%! Mean, Variance, and Covariance — Unit 8, Lesson 8.4 — Homework
\documentclass[10pt]{article}
\usepackage{linalg-article}
\usepackage{linalg-boxes}
\newcommand{\TallMath}[1]{$\displaystyle #1\rule[-1.4em]{0pt}{3.2em}$}

\begin{document}

\pageheader{Unit 8, Lesson 8.4}{Homework}
\namedateperiod

\vspace{0.12in}

\begin{notesbox}{Practice}
The \textbf{shape} of a data set is captured by its \textbf{mean} $\mu=\frac1N\sum x_i$ (center), its
\textbf{variance} $\sigma^2=\frac1N\sum(x_i-\mu)^2$ (spread), and, for two features, the \textbf{covariance}
$\sigma_{xy}=\frac1N\sum(x_i-\mu_x)(y_i-\mu_y)$ (how they move together). These pack into the symmetric
\textbf{covariance matrix} $V=\left[\begin{smallmatrix}\sigma_x^2 & \sigma_{xy}\\ \sigma_{xy} & \sigma_y^2\end{smallmatrix}\right]$.

\begin{enumerate}[leftmargin=1.9em, itemsep=11pt]
  \item \textbf{Spread, same center.} Two classes both average $7$ on a quiz.
    \begin{enumerate}[label=(\alph*), leftmargin=1.8em, itemsep=4pt]
      \item Class A: $4,6,8,10$. \hfill $\mu=\underline{\hspace{1.4cm}}$, \; $\sigma^2=\underline{\hspace{1.8cm}}$
      \item Class B: $1,5,9,13$. \hfill $\mu=\underline{\hspace{1.4cm}}$, \; $\sigma^2=\underline{\hspace{1.8cm}}$
    \end{enumerate}
    Same mean --- which class is more spread out, and how does the variance show it? \writeline
  \item \textbf{Covariance of two features.} Four data points $(x,y)$: $(1,2),(3,6),(5,4),(7,8)$.
        Find $\mu_x,\mu_y$, then the centered deviations, then the covariance $\sigma_{xy}$. Is it $+$, $-$, or $0$,
        and what does that mean? \writelines{3}
  \item \textbf{Build $V$.} The two features in \#2 each have variance $5$. Write the covariance matrix $V$.
        Check that it is symmetric. \writelines{2}
  \item \textbf{The PCA connection.} Your $V$ from \#3 is $\left[\begin{smallmatrix}5&4\\4&5\end{smallmatrix}\right]$.
        Find its eigenvalues (solve $(5-\lambda)^2-16=0$) and eigenvectors. Which eigenvector is PC1, and what
        fraction of the total variance does it capture? \writelines{3}
  \item \textbf{Explain.} (a) Why do we \emph{square} the distances from the mean for variance? (b) What does a
        \emph{negative} covariance say about two features? (c) In one phrase, why must the covariance matrix be
        symmetric? \writelines{3}
\end{enumerate}
\end{notesbox}

\begin{extensionbox}
\textbf{Tying it all together.}
\begin{enumerate}[label=(\alph*), leftmargin=1.8em, itemsep=5pt]
  \item \textbf{The matrix behind it.} For the notes' data (centered pairs $(-3,-1),(-1,-3),(1,3),(3,1)$), form
        the $2\times2$ matrix $A^{\top}\!A$ of dot products of the centered columns, then divide by $N=4$. What
        do you get, and why is ``the covariance matrix is $\frac1N A^{\top}\!A$'' the same PCA setup as Lesson~7.3?
  \item \textbf{Expected value.} A spinner lands on $1$ with probability $\tfrac23$ and on $4$ with probability
        $\tfrac13$. The \emph{expected value} is $E[x]=\sum p_i x_i$. Compute it. Why is it \emph{not} the
        midpoint $2.5$?
  \item \textbf{The whole course.} In two or three sentences, trace the arc: how does raw \emph{data} become a
        \emph{symmetric matrix} (Unit~6) whose \emph{eigen-directions} (Unit~7) reveal its shape --- and where do
        the mean, variance, and covariance enter?
\end{enumerate}
\writelines{5}
\end{extensionbox}

\begin{spiralbox}
\textbf{You've reached the end of the course.} From ``what is a vector'' (Unit~1) through solving systems,
subspaces, orthogonality and least squares (Unit~4), determinants, eigenvalues (Unit~6), the SVD and PCA
(Unit~7), and learning from data (Unit~8), one idea recurs: a matrix is a \emph{transformation}, and its special
directions tell you what it does. Mean, variance, and covariance are how real data walks in the door --- and the
covariance matrix hands it straight to everything you have learned. \emph{Well done.}
\end{spiralbox}

\end{document}
